% !TEX root = ../../main.tex
%% Défi du chapitre 12 — L'opérateur de différence et la formule de Newton
%% Résultats vérifiés par verifications/defis_11_15.py.

\rubrique{Défi}

\begin{defi}{L'opérateur de différence et la formule de Newton}{1 h 30}
Soit $n\in\N^{*}$ et $E=\R_n[X]$. On définit l'opérateur de différence
$\Delta:P\mapsto P(X+1)-P(X)$, version « discrète » de la dérivation, et
les polynômes
\[
H_0=1,\qquad H_k=\frac{X(X-1)\cdots(X-k+1)}{k!}\quad(k\geqslant1),
\]
qui vérifient $H_k(m)=\dbinom mk$ pour tout entier $m\geqslant0$.

\partieExo{A}{Un endomorphisme nilpotent}
\begin{questions}
  \item Montrer que $\Delta$ est un endomorphisme de $E$, et que si
    $\deg P=d\geqslant1$, alors $\deg\Delta P=d-1$ : préciser le
    coefficient dominant de $\Delta P$.
  \item Déterminer $\operatorname{Ker}\Delta$ et $\operatorname{Im}\Delta$
    (théorème du rang). Montrer que $\Delta^{n+1}=0$ mais
    $\Delta^{n}\neq0$.
  \item Écrire la matrice de $\Delta$ dans la base canonique de $\R_3[X]$.
\end{questions}

\partieExo{B}{La base binomiale et la formule de Newton}
\begin{questions}[resume]
  \item Montrer que $(H_0,\dots,H_n)$ est une base de $E$, et que
    $\Delta H_k=H_{k-1}$ pour $k\geqslant1$ : c'est la formule de Pascal
    pour les polynômes. Écrire la matrice de $\Delta$ dans cette base.
  \item En déduire la \emph{formule de Newton} : pour tout $P\in E$,
    \[
    P=\sum_{k=0}^{n}\bigl(\Delta^{k}P\bigr)(0)\,H_k .
    \]
    Décomposer ainsi $X^{3}$, en lisant les $\bigl(\Delta^{k}P\bigr)(0)$
    sur la table des différences ci-dessous.
  \item Montrer que $P$ prend des valeurs entières en tout entier relatif
    si et seulement si ses coordonnées dans la base $(H_k)$ sont entières.
    Exemple : $\frac{X(X+1)}{2}$, à coefficients non entiers, mais à
    valeurs entières.
\end{questions}

\medskip
\begin{center}
\begin{tikzpicture}[x=1.05cm,y=0.72cm]
  \foreach \v [count=\i from 0] in {0,1,8,27,64,125}{%
    \node[font=\small,text=otInk] at (\i,0) {$\v$};}
  \foreach \v [count=\i from 0] in {1,7,19,37,61}{%
    \node[font=\small,text=otNavy] at (\i+0.5,-1) {$\v$};}
  \foreach \v [count=\i from 0] in {6,12,18,24}{%
    \node[font=\small,text=otNavy] at (\i+1,-2) {$\v$};}
  \foreach \v [count=\i from 0] in {6,6,6}{%
    \node[font=\small,text=otNavy] at (\i+1.5,-3) {$\v$};}
  \foreach \v [count=\i from 0] in {0,0}{%
    \node[font=\small,text=otMuted] at (\i+2,-4) {$\v$};}
  \foreach \xx/\yy in {0/0,0.5/-1,1/-2,1.5/-3}{%
    \draw[otGreenDark,line width=0.8pt,rounded corners=2pt] (\xx-0.27,\yy-0.3) rectangle (\xx+0.27,\yy+0.3);}
  \foreach \l/\yy in {$P(k)=k^{3}$/0,$\Delta P$/-1,$\Delta^{2}P$/-2,$\Delta^{3}P$/-3,$\Delta^{4}P$/-4}{%
    \node[font=\scriptsize,text=otMuted,anchor=east] at (-0.6,\yy) {\l};}
\end{tikzpicture}\\[0.5mm]
\legende{La table des différences de $k\mapsto k^{3}$ : chaque ligne est
formée des différences de la précédente ; la diagonale encadrée donne
$\bigl(\Delta^{k}P\bigr)(0)=0,1,6,6$.}
\end{center}

\partieExo{C}{Sommation discrète}
\begin{questions}[resume]
  \item Montrer que si $\Delta Q=P$, alors
    $\sum_{k=0}^{m-1}P(k)=Q(m)-Q(0)$ pour tout $m\in\N^{*}$ :
    $Q$ joue le rôle d'une « primitive discrète » de $P$.
  \item Trouver $Q$ tel que $\Delta Q=X^{2}$, et en déduire
    $\sum_{k=0}^{m-1}k^{2}$. Faire de même pour
    $\sum_{k=0}^{m-1}k^{3}$, et retrouver le théorème de
    Nicomaque (défi du chapitre 2).
\end{questions}
\end{defi}
